\section{Bounded potentials and recursive Hessian currents}
\label{sec:hessian-bounded-currents}\index{Hessian measure!bounded recursive currents}

The decreasing approximation of bounded admissible potentials supplies
limits that uniform approximation cannot reach. A bounded potential can
multiply a current with measure coefficients, so one can apply $\ddc$
to the resulting current. Repeating this operation requires bounds on
every coefficient, rather than only on the trace of a positive mixed
measure. The absolute-form estimate provides these bounds below degree
$m$; at degree $m$, positivity is obtained after wedging with
$\omega^{n-m}$.

All bounded potentials have their canonical upper semicontinuous representatives from
\eqref{eq:hessian-bounded-upper-tests}. Their pointwise values will be
used when multiplying measure coefficients. For a real current $T$ of
order zero, write $\|T\|_K$ for the sum of the total variations of its
coefficients on a compact coordinate set $K$. Fix $K\Subset U$ in a
coordinate chart. Different fixed coordinate frames give equivalent
local bounds.

\Needspace{11\baselineskip}
\begin{lemma}\label{hessian:bounded-zero-order}
If $u$ is bounded admissible, then $\ddc u$ has signed Radon measure
coefficients. For $K\Subset U$ contained in a coordinate chart,
\[
 \|\ddc u\|_K\leq C_{K,U,\omega}(1+\|u\|_{L^\infty(X)}).
\]
\end{lemma}
\begin{proof}
Choose smooth strictly admissible $u_j\downarrow u$ by
\cref{hessian:bounded-decreasing-approximation}, and put
$\sigma_j=\omega+\ddc u_j$. After discarding finitely many terms,
$\|u_j\|_\infty\leq\|u\|_\infty+1$. Indeed, the increasing open sets
$\{u_j<\sup_Xu+1\}$ cover $X$, and compactness gives a single index
whose open set is all of $X$; the lower bound follows from $u_j\geq u$.

The closed $m$-cone lies in the closed $2$-cone. In an $\omega$-unitary
frame, its eigenvalues satisfy
\[
 \sum_i\lambda_i^2\leq(\operatorname{tr}\sigma_j)^2,
 \qquad \operatorname{tr}\sigma_j\geq0.
\]
Thus every coefficient of $\sigma_j$ is bounded by its trace. The
degree-one Chern--Levine--Nirenberg estimate in
\cref{hessian:continuous-products}, applied to these smooth functions,
bounds the integral of that trace on $K$ by
$C(1+\|u\|_\infty)$. Subtracting the fixed smooth $\omega$ gives the
same bound for the coefficient variations of $\ddc u_j$.
Bounded convergence gives $u_j\to u$ in local $L^1$, and therefore
$\ddc u_j\to\ddc u$ in distributions. The uniform coefficient bounds
show that this limit has measure coefficients with the asserted bound.
\end{proof}

Set $D_0=1$. For an ordered list $u_1,\ldots,u_p$ of bounded
admissible potentials, define recursively
\begin{equation}\label{eq:hessian-bounded-recursion}
 D_p(u_1,\ldots,u_p)
 =\ddc\bigl(u_pD_{p-1}(u_1,\ldots,u_{p-1})\bigr).
\end{equation}
The next result justifies the bounded multiplication at every step up
to $p=m$.

\Needspace{13\baselineskip}
\begin{proposition}\label{hessian:bounded-recursive-currents}
For $1\leq p<m$, the recursion defines a real closed current of order
zero and bidegree $(p,p)$. If
$M=\max_a\|u_a\|_{L^\infty(X)}$, then
\[
 \|D_p(u_1,\ldots,u_p)\|_K
       \leq C_{K,U,\omega,p}(1+M)^p.
\]
For $p=m$, it defines a real closed distributional current. For any
smooth admissible sequences $u_a^{j_a}\downarrow u_a$ with a common
uniform bound, one has
\begin{equation}\label{eq:hessian-bounded-iterated-limit}
 D_p(u_1,\ldots,u_p)
 =\lim_{j_p\to\infty}\cdots\lim_{j_1\to\infty}
    \ddc u_1^{j_1}\wedge\cdots\wedge\ddc u_p^{j_p}
\end{equation}
in distributions. The index $j_1$ tends to infinity first. Below degree
$m$, all these limits have locally uniformly bounded coefficient
variations.
\end{proposition}
\begin{proof}
The case $p=1$ is the preceding lemma. Suppose the assertions hold
through degree $p-1<m$. Since $u_p$ is bounded Borel and $D_{p-1}$
has measure coefficients, their product is well defined. Applying
$\ddc$ gives a real current that is both $\mathrm d$- and
$\dc$-closed.

For a fixed smooth complementary test form $\eta$, dominated
convergence against the fixed signed measures gives
\[
 \begin{aligned}
 \langle D_p,\eta\rangle
 &=\langle u_pD_{p-1},\ddc\eta\rangle\\
 &=\lim_{j_p\to\infty}
       \langle u_p^{j_p}D_{p-1},\ddc\eta\rangle.
 \end{aligned}
\]
With $j_p$ fixed, the inductive iterated limit for $D_{p-1}$ can be
paired with the smooth form $u_p^{j_p}\ddc\eta$. Integration by parts
for the smooth closed factors identifies the resulting pairing with
the smooth product on the right of
\eqref{eq:hessian-bounded-iterated-limit}. This proves that identity
also when $p=m$.

For $p<m$, expand the smooth product using
$\ddc u_a^{j_a}=\sigma_a^{j_a}-\omega$. A term with $q$ admissible
factors, $0\leq q\leq p$, has the form
\[
 \sigma_{a_1}^{j_{a_1}}\wedge\cdots\wedge
 \sigma_{a_q}^{j_{a_q}}\wedge\omega^{p-q}.
\]
Every such $q$ is strictly less than $m$. Apply
\cref{hessian:absolute-mixed-forms} to each coefficient test, and then
the smooth mass estimate from \cref{hessian:continuous-products} to
the potential sum. The term $q=0$ is the fixed smooth form $\omega^p$.
For $q\geq1$, the sum of its admissible forms is
$q(\omega+\ddc\bar u)$, where $\bar u$ is their average and
$\|\bar u\|_\infty\leq M+1$. The result is a
coefficient variation bound $C(1+M)^p$, after choosing the
approximations with norm at most $M+1$ as in the preceding proof.
For arbitrary uniformly bounded approximations the same argument
uses their common bound. The bounds survive each distributional
limit, proving that $D_p$ has order zero and the stated estimate.
This completes the induction. No coefficient bound for the unwedged
degree-$m$ current is used or asserted.
\end{proof}

\Needspace{8\baselineskip}
For a subset $I=\{i_1<\cdots<i_q\}\subset\{1,\ldots,p\}$, write
$D_I=D_q(u_{i_1},\ldots,u_{i_q})$, with $D_\varnothing=1$.
Define the shifted current and its mixed positive current by
\begin{equation}\label{eq:hessian-bounded-shifted-current}
 R_p(u_1,\ldots,u_p)
 =\sum_{I\subset\{1,\ldots,p\}}\omega^{p-|I|}\wedge D_I,
 \qquad Q_p=R_p\wedge\omega^{n-m}.
\end{equation}
Below degree $m$, $R_p$ has measure coefficients. At degree $m$ it
is initially a distributional current. Its positive wedge $Q_m$ is
nevertheless a measure.

Put $M=\max_a\|u_a\|_{L^\infty(X)}$.

\Needspace{13\baselineskip}
\begin{theorem}\label{hessian:bounded-shifted-positive}\index{Hessian measure!iterated decreasing limits}
For $1\leq p\leq m$, $Q_p$ is a positive current of order zero, and
\[
 \int_KQ_p\wedge\omega^{m-p}
       \leq C_{K,U,\omega,p}(1+M)^p.
\]
It is the iterated limit, in the order of
\eqref{eq:hessian-bounded-iterated-limit}, of
\[
 (\omega+\ddc u_1^{j_1})\wedge\cdots\wedge
 (\omega+\ddc u_p^{j_p})\wedge\omega^{n-m}.
\]
In particular $Q_m$ is a positive Radon measure. These currents are
local in the canonical potentials and independent of the chosen
approximations for the fixed ordering. When all potentials are
continuous, $Q_p$ equals the continuous mixed product of
\cref{hessian:continuous-products}.
\end{theorem}
\begin{proof}
Expand the smooth shifted product over its subsets. Applying
\eqref{eq:hessian-bounded-iterated-limit} to each subset, in its
inherited ordering, gives \eqref{eq:hessian-bounded-shifted-current}.
The absent indices have no effect on a subset's limit.

Each smooth product is positive: pairing with a decomposable
strongly positive complementary test appends $m-p$ nonnegative
rank-one $(1,1)$-forms, and polarized G\aa rding positivity applies.
The smooth mixed mass bound is uniform in all indices. Distributional
limits preserve positivity and the local trace bound at every step.
Positive currents have order zero, so this gives the assertions for
$Q_p$, including the top-degree measure.

The recursion uses only bounded multiplication of canonical values
and local distributional differentiation. It is therefore local.
Its definition contains no approximation choices, so the proved
limit is independent of those choices while the ordering is fixed.
For continuous potentials, the approximations converge uniformly by
Dini's theorem. The telescoping estimate in
\cref{hessian:continuous-products} is uniform in all remaining
indices and shows that their iterated limit is the same continuous
mixed product. This also proves permutation invariance in the
continuous case.
\end{proof}

The shifted currents are not closed. Their derivatives can be stated
without multiplying a discontinuous function by a distribution of
positive order. This is the torsion identity needed before considering
integration by parts with bounded weights.

\Needspace{15\baselineskip}
\begin{proposition}\label{hessian:bounded-current-torsion}\index{Torsion!bounded Hessian currents}
Write $R_{p-1}^{\widehat a}$ for the shifted current obtained by
omitting $u_a$, in the inherited order, and use two hats for two
omissions. Put $R_0=1$. For $1\leq p\leq m$,
\[
 \begin{aligned}
 \mathrm dR_p&=\mathrm d\omega\wedge
                         \sum_{a=1}^pR_{p-1}^{\widehat a},\\
 \dc R_p&=\dc\omega\wedge
                         \sum_{a=1}^pR_{p-1}^{\widehat a},\\
 \ddc R_p&=\ddc\omega\wedge
                         \sum_{a=1}^pR_{p-1}^{\widehat a}
       +2\,\mathrm d\omega\wedge\dc\omega\wedge
                         \sum_{a<b}R_{p-2}^{\widehat a\widehat b}.
 \end{aligned}
\]
The last sum is absent for $p=1$. All currents on the right have order
zero. If all but at most one potential are smooth, $R_p$ is independent
of the ordering and is the ordinary
product of the smooth factors with $\omega+\ddc u$.
\end{proposition}
\begin{proof}
Every $D_I$ is both $\mathrm d$- and $\dc$-closed. Differentiating
\eqref{eq:hessian-bounded-shifted-current} therefore gives
\[
 \mathrm dR_p=\mathrm d\omega\wedge
 \sum_I(p-|I|)\omega^{p-|I|-1}\wedge D_I,
\]
with zero terms omitted. The coefficient $p-|I|$ counts the omitted
indices outside $I$, which proves the first formula; the second is
identical with $\dc$ in place of $\mathrm d$. Apply $\mathrm d$ to
the second formula. Each pair of omitted indices occurs twice, and
$-\dc\omega\wedge\mathrm d\omega
=\mathrm d\omega\wedge\dc\omega$ because both forms have odd
degree three. This gives the displayed $\ddc$ formula, with every
background derivative retained.

If all but one potential are smooth, applying $\ddc$ to the bounded
potential times a smooth closed product gives its $\ddc$ wedged
with that product. At later smooth steps, closedness again makes
the recursion ordinary smooth multiplication. This is independent
of the position of the bounded potential, and the subset expansion
gives the last assertion.
\end{proof}

The iterated limit fixes the order of the potentials. To compare it
with a simultaneous limit, discontinuous weights must be controlled
against varying coefficient measures. The next section supplies this
control for the first Hessian, giving symmetry and simultaneous
decreasing convergence for two factors in
\cref{hessian:bounded-two-factor-convergence}. Higher-degree products
still need additional estimates.
